\documentclass[12pt,a4paper]{article}
\usepackage{iftex}
\ifPDFTeX
  \usepackage[T2A]{fontenc}
  \usepackage[utf8]{inputenc}
  \usepackage[english,russian]{babel}
\else
  \usepackage{fontspec}
  \setmainfont{DejaVu Serif}
  \setmonofont{DejaVu Sans Mono}
  \renewcommand{\contentsname}{Содержание}
  \renewcommand{\figurename}{Рис.}
  \renewcommand{\tablename}{Таблица}
\fi

\usepackage[a4paper,margin=25mm]{geometry}
\usepackage{amsmath,amssymb,amsthm}
\usepackage{graphicx}
\usepackage{booktabs}
\usepackage{enumitem}
\usepackage[hidelinks]{hyperref}

\newcommand{\Exp}{\mathsf{M}}
\newcommand{\Var}{\mathsf{D}}
\newcommand{\PP}{\mathsf{P}}
\newcommand{\dens}{\mathrm{p}}
\newcommand{\Normal}{\mathcal{N}}
\newcommand{\R}{\mathbb{R}}
\newcommand{\xs}[1]{x^{(#1)}}
\newcommand{\ys}[1]{y^{(#1)}}
\newcommand{\norm}[1]{\left\lVert #1 \right\rVert}
\newcommand{\abs}[1]{\left\lvert #1 \right\rvert}
\newcommand{\dd}{\,\mathrm{d}}
\newcommand{\est}[1]{\tilde{#1}}
\newcommand{\zad}[1]{\medskip\noindent\textbf{Задание #1.}\quad}
\newcommand{\prov}{\smallskip\noindent\textit{Проверка:}\ }

\theoremstyle{definition}
\newtheorem{definition}{Определение}
\newtheorem{statement}{Утверждение}
\theoremstyle{plain}
\newtheorem{lemma}{Лемма}

% Список с русскими буквенными метками
\newenvironment{ruslist}
  {\begin{itemize}[label={},leftmargin=2.2em,itemsep=2pt,topsep=3pt]}
  {\end{itemize}}


\begin{document}

\begin{center}
  {\large\bfseries Задача 3. Метод наименьших квадратов}\\[1mm]
  Нейронные сети и машинное обучение\\[1mm]
  Задания 1--2
\end{center}
\vspace{2mm}

Дана выборка $\{(\xs{k}, \ys{k})\}_{k=1}^{N}$, ищется функция
$\est f(x, w)$, дающая минимум
\[
  E(w) = \frac12\sum_{k=1}^{N}\bigl(\ys{k} - \est f(\xs{k}, w)\bigr)^2 .
\]
Ниже $x_k = \xs{k}$, $y_k = \ys{k}$. Числа в проверках получены скриптом
\texttt{check.py} с фиксированным зерном.

\zad{1} Линейная регрессия $\est f(x) = Ax + B$. Условие минимума
$\partial E/\partial A = \partial E/\partial B = 0$ даёт систему нормальных
уравнений
\begin{equation}\label{eq:lin}
  \begin{cases}
    \bigl(\sum_k x_k^2\bigr)A + \bigl(\sum_k x_k\bigr)B = \sum_k x_k y_k,\\[1mm]
    \bigl(\sum_k x_k\bigr)A + N B = \sum_k y_k .
  \end{cases}
\end{equation}
Указать, когда \eqref{eq:lin} разрешима, почему решение единственно, и какой
выборке отвечает нулевой определитель.

\emph{Определитель.} Матрица системы
\[
  G = \begin{pmatrix} \sum_k x_k^2 & \sum_k x_k \\ \sum_k x_k & N \end{pmatrix}
    = \Phi^{\top}\Phi, \qquad
  \Phi = \begin{pmatrix} x_1 & 1 \\ \vdots & \vdots \\ x_N & 1 \end{pmatrix},
\]
то есть $G$ -- матрица Грама столбцов $\Phi$. Обозначив
$\bar x = \tfrac1N\sum_k x_k$,
\[
  \det G = N\sum_k x_k^2 - \Bigl(\sum_k x_k\Bigr)^{\!2}
         = N\sum_k x_k^2 - N^2\bar x^2
         = N\sum_k (x_k - \bar x)^2 \ge 0 .
\]
Значит
\[
  \det G = 0 \iff x_k = \bar x \ \ \forall k \iff
  \text{все } x_k \text{ совпадают},
\]
и система однозначно разрешима тогда и только тогда, когда среди $x_k$ есть хотя
бы два различных.

\emph{Единственность.} $E(A,B)$ -- квадратичная функция, её матрица вторых
производных равна $G$. При $\det G > 0$ и $\mathrm{tr}\,G > 0$ матрица $G$
положительно определена, значит $E$ строго выпукла, стационарная точка
единственна и является точкой глобального минимума. При $\det G = 0$ строгой
выпуклости нет.

\emph{Вырожденный случай.} $\det G = 0$ означает, что все измерения выполнены в
одной точке $x_k \equiv x_0$: точки выборки лежат на вертикальной прямой. Тогда
\[
  E(A,B) = \frac12\sum_k \bigl(y_k - (Ax_0 + B)\bigr)^2
\]
зависит только от комбинации $t = Ax_0 + B$ и минимальна при $t = \bar y$.
Решением служит вся прямая $Ax_0 + B = \bar y$ в плоскости $(A,B)$: наклон $A$
данными не определён, ранг $G$ равен 1.

\prov $N = 12$, $x_k$ равномерны на $[-2,2]$, $y_k = 3x_k - 1 + \Normal(0,0.3^2)$:
$\det G = 129.7672$ совпадает с $N\sum_k(x_k-\bar x)^2 = 129.7672$, решение
$A = 3.1677$, $B = -1.1930$.

Для $x_k \equiv 1.7$: $\det G = 1.2\cdot10^{-13}$, ранг 1, и на прямой
$Ax_0 + B = \bar y$ значение ошибки одно и то же:
$E(0,\,4.030) = E(3,\,-1.070) = E(-7.5,\,16.780) = 0.295260$.

\zad{2} Полиномиальная регрессия
$\est f(x, w) = \sum_{j=0}^{M} w_j x^j$. Показать, что минимизирующие $E(w)$
коэффициенты удовлетворяют системе
\begin{equation}\label{eq:poly}
  \sum_{j=0}^{M} a_{ij} w_j = b_i, \qquad
  a_{ij} = \sum_{k=1}^{N} x_k^{\,i+j}, \qquad
  b_i = \sum_{k=1}^{N} x_k^{\,i} y_k, \qquad i = 0,\dots,M,
\end{equation}
и выяснить, всегда ли она разрешима.

\emph{Вывод системы.}
\[
  \frac{\partial E}{\partial w_i}
  = -\sum_{k=1}^{N}\Bigl(y_k - \sum_{j=0}^{M} w_j x_k^{\,j}\Bigr) x_k^{\,i} = 0
  \;\Longrightarrow\;
  \sum_{j=0}^{M}\Bigl(\underbrace{\sum_{k} x_k^{\,i+j}}_{a_{ij}}\Bigr) w_j
  = \underbrace{\sum_{k} x_k^{\,i} y_k}_{b_i}. \qquad\blacksquare
\]

\emph{Разрешимость.} В матричном виде $A w = b$, где
\[
  A = V^{\top}V, \quad b = V^{\top}y, \quad
  V = \bigl(x_k^{\,j}\bigr)_{k=1..N,\ j=0..M}
\]
-- матрица Вандермонда размера $N \times (M+1)$. Как матрица Грама, $A$
симметрична и неотрицательно определена.

Ядро: $Aw = 0 \Rightarrow w^{\top}V^{\top}Vw = \norm{Vw}^2 = 0 \Rightarrow Vw = 0$,
то есть $\ker A = \ker V$ и $\mathrm{rank}\,A = \mathrm{rank}\,V$. Условие
$Vw = 0$ означает, что многочлен $\sum_j w_j x^j$ степени не выше $M$ обращается
в нуль во всех узлах $x_k$. Ненулевой многочлен степени $\le M$ имеет не более
$M$ корней, поэтому, обозначив через $r$ число \emph{различных} узлов,
\[
  \mathrm{rank}\,A = \min(r,\, M+1), \qquad
  \det A \ne 0 \iff r \ge M+1 .
\]

Совместность выполняется всегда: из $\ker A = \ker V$ следует
$\mathrm{rank}\,A = \mathrm{rank}\,V^{\top}$, а вместе с
$\mathrm{Im}\,A \subseteq \mathrm{Im}\,V^{\top}$ это даёт
$\mathrm{Im}\,A = \mathrm{Im}\,V^{\top}$; правая часть
$b = V^{\top}y$ лежит в $\mathrm{Im}\,V^{\top}$. Итого:
\begin{ruslist}
  \item[--] $r \ge M+1$: решение существует и единственно;
  \item[--] $r \le M$: решение существует, но их бесконечно много -- аффинное
    подпространство размерности $M+1-r$; все они дают одинаковые значения
    многочлена в узлах и одинаковую ошибку $E$.
\end{ruslist}

Формальной невырожденности мало. При $x \in [0,1]$
$a_{ij} = \sum_k x_k^{i+j} \approx N/(i+j+1)$, то есть $A$ близка к матрице
Гильберта, умноженной на $N$, и её число обусловленности растёт с $M$
экспоненциально. Начиная примерно с $M = 10$ матрица теряет ранг численно, хотя
формально невырождена.

\prov $N = 50$ различных узлов, равномерных на $[0,1]$.

\begin{center}
\begin{tabular}{ccc}
\toprule
$M$ & $\mathrm{cond}\,A$ & численный ранг из $M+1$ \\
\midrule
$0$  & $1.0$              & $1$ из $1$ \\
$1$  & $2.0\cdot10^{1}$   & $2$ из $2$ \\
$3$  & $1.7\cdot10^{4}$   & $4$ из $4$ \\
$5$  & $1.4\cdot10^{7}$   & $6$ из $6$ \\
$10$ & $1.2\cdot10^{15}$  & $10$ из $11$ \\
$15$ & $2.8\cdot10^{17}$  & $12$ из $16$ \\
\bottomrule
\end{tabular}
\end{center}

Выборка из $r = 4$ различных узлов, каждый повторён 5 раз ($N = 20$):

\begin{center}
\begin{tabular}{cccc}
\toprule
$M$ & $\mathrm{rank}\,A$ из $M+1$ & $\dim\ker A$ & $\norm{Aw - b}$ \\
\midrule
$3$ & $4$ из $4$ & $0$ & $1.6\cdot10^{-14}$ \\
$5$ & $4$ из $6$ & $2$ & $8.3\cdot10^{-15}$ \\
$8$ & $4$ из $9$ & $5$ & $9.1\cdot10^{-15}$ \\
\bottomrule
\end{tabular}
\end{center}

Ранг равен $r$ независимо от $M$, размерность ядра равна $M+1-r$, невязка нулевая
с точностью до округления: система совместна и при вырожденной $A$.

\end{document}
